\documentclass[11pt,a4paper]{article}

% ─── Paquetes Base y Codificación ───────────────────────────────────────────
\usepackage[utf8]{inputenc}
\usepackage[T1]{fontenc}
\usepackage[spanish,es-tabla]{babel}
\usepackage{geometry}
\geometry{margin=2.5cm}

% ─── Tipografía y Diseño Visual ─────────────────────────────────────────────
\usepackage{lmodern}
\usepackage{microtype}
\usepackage{xcolor}
\usepackage{titlesec}
\usepackage{enumitem}
\usepackage{booktabs}
\usepackage{tabularx}
\usepackage{amsmath,amssymb}
\usepackage{listings}
\usepackage{hyperref}
\usepackage{fancyhdr}

% ─── Paleta de Colores Corporativos Slotify ─────────────────────────────────
\definecolor{SlotifyIndigo}{HTML}{6366F1}
\definecolor{SlotifyViolet}{HTML}{8B5CF6}
\definecolor{SlotifyDark}{HTML}{171717}
\definecolor{SlotifyGray}{HTML}{525252}
\definecolor{SlotifyLightBg}{HTML}{F8FAFC}
\definecolor{CodeBorder}{HTML}{E2E8F0}

\hypersetup{
    colorlinks=true,
    linkcolor=SlotifyIndigo,
    urlcolor=SlotifyViolet,
    pdftitle={Justificación Arquitectónica SOLID - Identidad de Marca en Formularios (Slotify)},
    pdfauthor={Equipo de Desarrollo Slotify}
}

% ─── Configuración de Bloques de Código ─────────────────────────────────────
\lstdefinestyle{slotifystyle}{
    backgroundcolor=\color{SlotifyLightBg},
    commentstyle=\color{SlotifyGray}\itshape,
    keywordstyle=\color{SlotifyIndigo}\bfseries,
    stringstyle=\color{SlotifyViolet},
    basicstyle=\ttfamily\footnotesize,
    breakatwhitespace=false,
    breaklines=true,
    captionpos=b,
    keepspaces=true,
    numbers=left,
    numbersep=8pt,
    numberstyle=\tiny\color{SlotifyGray},
    frame=single,
    rulecolor=\color{CodeBorder},
    showspaces=false,
    showstringspaces=false,
    showtabs=false,
    tabsize=2
}
\lstset{style=slotifystyle}

% ─── Encabezado y Pie de Página ─────────────────────────────────────────────
\pagestyle{fancy}
\fancyhf{}
\lhead{\textbf{Slotify} $\cdot$ Arquitectura de Software y Principios SOLID}
\rhead{Identidad Corporativa en Formularios}
\cfoot{\thepage}

% ─── Formato de Títulos ─────────────────────────────────────────────────────
\titleformat{\section}
  {\normalfont\Large\bfseries\color{SlotifyIndigo}}{\thesection.}{0.6em}{}
\titleformat{\subsection}
  {\normalfont\large\bfseries\color{SlotifyDark}}{\thesubsection}{0.5em}{}

% ─── Documento ──────────────────────────────────────────────────────────────
\begin{document}

\begin{titlepage}
    \centering
    \vspace*{2cm}
    {\scshape\LARGE Plataforma Slotify / Slotly \par}
    \vspace{0.8cm}
    {\huge\bfseries Fundamentación Arquitectónica y Cumplimiento de Principios SOLID\par}
    \vspace{0.6cm}
    {\Large Historia de Usuario: Carga de Logotipo y Colores Corporativos Principales sobre Formularios Públicos de Reserva\par}
    \vspace{1.5cm}

    \begin{tabular}{rl}
        \toprule
        \textbf{Rol:} & Administrador de un negocio en Slotly \\
        \textbf{Objetivo:} & Cargar logotipo y definir colores corporativos principales \\
        \textbf{Propósito:} & Mantener la identidad de marca en la interfaz pública de reservas \\
        \textbf{Arquitectura Frontend:} & Feature-Based Architecture (FBA) + MVVM Reactivo (Expo / React Native) \\
        \textbf{Arquitectura Backend:} & Clean Architecture + CQRS con MediatR y FluentValidation (.NET 9) \\
        \bottomrule
    \end{tabular}

    \vfill
    {\large Fecha de generación: \today\par}
\end{titlepage}

\tableofcontents
\newpage

% ============================================================================
\section{Introducción y Regla de Dependencia Concéntrica}
% ============================================================================

La funcionalidad de personalización de identidad de marca (logotipo institucional, color corporativo primario \texttt{accentColor} y color corporativo secundario \texttt{secondaryColor}) sobre los formularios públicos de reserva fue diseñada siguiendo la \textbf{Regla de Dependencia Concéntrica} de \textit{Clean Architecture} en el backend (\texttt{Slotify\_back}) y \textit{Feature-Based Architecture (FBA)} con el patrón \textbf{MVVM} en el frontend (\texttt{Slotify}).

Formalmente, el flujo de dependencias entre capas apunta estrictamente hacia el núcleo del dominio:

\[
\text{Frameworks / UI / API REST} \longrightarrow \text{Casos de Uso / Application (CQRS)} \longrightarrow \text{Dominio / Entidades}
\]

En el backend (.NET 9), esto se expresa como:
\[
\underbrace{\texttt{Slotify.API} \cup \texttt{Slotify.Infrastructure}}_{\text{Capas Externas (Detalles)}}
\longrightarrow
\underbrace{\texttt{Slotify.Application}}_{\text{Orquestación CQRS}}
\longrightarrow
\underbrace{\texttt{Slotify.Domain}}_{\text{Núcleo Invariante}}
\]

A continuación, se demuestra técnicamente cómo cada componente satisface los cinco principios \textbf{SOLID}.

% ============================================================================
\section{S --- Single Responsibility Principle (Responsabilidad Única)}
% ============================================================================

\begin{quote}
\textit{``Todo módulo, clase o función debe tener una única responsabilidad y una sola razón para cambiar.''}
\end{quote}

La solución evita el anti-patrón \textit{God Object} dividiendo las responsabilidades en unidades cohesivas tanto en el cliente móvil/web como en el servidor:

\subsection{Aplicación en el Backend (.NET 9)}
\begin{enumerate}[label=\textbf{\arabic*.}]
    \item \textbf{Entidad de Dominio (\texttt{Business.cs}):} Concentra únicamente el estado del negocio y las reglas invariantes de validación de identidad corporativa (\texttt{IsValidCorporateHexColor} e \texttt{IsValidLogoSource}), sin acoplarse a serializadores JSON ni a bases de datos.
    \item \textbf{Validadores de Entrada (\texttt{UpdateBusinessBrandIdentityValidator} y \texttt{UpdateBusinessFormConfigValidator}):} Su única razón de cambio es la modificación de las reglas sintácticas de entrada mediante \texttt{FluentValidation}.
    \item \textbf{Manejadores de Comandos (\texttt{UpdateBusinessBrandIdentityCommandHandler} y \texttt{UpdateBusinessFormConfigCommandHandler}):} Orquestan exclusivamente la persistencia transaccional de la identidad de marca y estructura del formulario.
    \item \textbf{Manejador de Consultas (\texttt{GetBusinessBookingConfigQueryHandler}):} Se encarga únicamente de recuperar y proyectar el esquema público del formulario (\texttt{BookingFormConfigDto}) a partir del \texttt{slug} del negocio.
\end{enumerate}

\subsection{Aplicación en el Frontend (React Native / TypeScript)}
\begin{enumerate}[label=\textbf{\arabic*.}]
    \item \textbf{Modelo y Persistencia (\texttt{formBranding.ts}):} Aísla las paletas corporativas predefinidas (\texttt{CORPORATE\_COLOR\_PRESETS}), la validación hexadecimal (\texttt{isValidHexColor}, \texttt{normalizeHexColor}) y la persistencia local (\texttt{saveFormSchemaLocally}).
    \item \textbf{Componente de Edición de Marca (\texttt{FormBrandIdentityCard.tsx}):} Vista pasiva encargada exclusivamente de capturar el logotipo y los colores corporativos, notificando cambios mediante la propiedad \texttt{onChangeBranding}.
    \item \textbf{Renderizador Público (\texttt{FormViewer.tsx}):} Se dedica únicamente a proyectar el formulario de reserva con el logotipo y los colores corporativos del establecimiento en los modos \texttt{classic\_scroll} e \texttt{interactive\_slides}.
\end{enumerate}

% ============================================================================
\section{O --- Open/Closed Principle (Principio de Abierto/Cerrado)}
% ============================================================================

\begin{quote}
\textit{``Las entidades de software deben estar abiertas para su extensión, pero cerradas para su modificación.''}
\end{quote}

\begin{itemize}
    \item \textbf{Extensibilidad del Pipeline de Validación (Backend):} Gracias al comportamiento transversal \texttt{ValidationBehavior<TRequest, TResponse>} registrado en MediatR, la incorporación de \texttt{UpdateBusinessBrandIdentityValidator} extiende automáticamente las validaciones del sistema sin modificar una sola línea del pipeline ni de los controladores existentes:
\begin{lstlisting}[language={[Sharp]C}]
public class ValidationBehavior<TRequest, TResponse>(
    IEnumerable<IValidator<TRequest>> validators)
    : IPipelineBehavior<TRequest, TResponse>
    where TRequest : notnull
\end{lstlisting}
    \item \textbf{Extensibilidad de Paletas Corporativas y Bloques de Formulario (Frontend):} La incorporación de nuevas combinaciones corporativas en \texttt{CORPORATE\_COLOR\_PRESETS} o nuevos tipos de preguntas en \texttt{FIELD\_TYPE\_META} extiende las capacidades visuales del constructor sin alterar la lógica interna de \texttt{FormBrandIdentityCard.tsx} ni el estado de \texttt{FormBuilderScreen.tsx}.
\end{itemize}

% ============================================================================
\section{L --- Liskov Substitution Principle (Sustitución de Liskov)}
% ============================================================================

\begin{quote}
\textit{``Si $S$ es un subtipo de $T$, entonces los objetos de tipo $T$ pueden ser sustituidos por objetos de tipo $S$ sin alterar las propiedades deseables del programa.''}
\end{quote}

Sea $T = \texttt{IBusinessRepository}$ el contrato abstracto definido en \texttt{Slotify.Domain} y sean $S_1 = \texttt{BusinessRepository}$ (implementación PostgreSQL con EF Core) y $S_2 = \texttt{Mock<IBusinessRepository>}$ (doble de prueba en \texttt{Slotify.Tests}). Para todo comando $c \in \texttt{UpdateBusinessBrandIdentityCommand}$:

\[
\forall r \in \{S_1, S_2\} \subseteq \texttt{IBusinessRepository}, \quad \texttt{Handler}(r).\texttt{Handle}(c) \implies \text{Contrato Postcondicional Invariante}
\]

\begin{itemize}
    \item \textbf{Sustitución Polimórfica de Repositorios:} Los manejadores CQRS operan de manera idéntica tanto con la persistencia real en Supabase/PostgreSQL (\texttt{BusinessRepository}) como con implementaciones en memoria durante la ejecución de pruebas unitarias (\texttt{BusinessBrandingTests.cs}), sin lanzar excepciones no contempladas en el contrato.
    \item \textbf{Compatibilidad de Contratos DTO (\texttt{BookingFormConfigDto}):} Los esquemas personalizados con logotipo y colores corporativos sustituyen de forma transparente a las configuraciones predeterminadas de plantillas (\texttt{SectorTemplate.FormConfig}), preservando valores por defecto seguros (\texttt{\#6366F1} y \texttt{\#8B5CF6}) cuando un negocio aún no ha cargado su logotipo.
\end{itemize}

% ============================================================================
\section{I --- Interface Segregation Principle (Segregación de Interfaces)}
% ============================================================================

\begin{quote}
\textit{``Ningún cliente debe verse forzado a depender de métodos o propiedades que no utiliza.''}
\end{quote}

En lugar de utilizar una interfaz monolítica de datos o propiedades globales acopladas, los contratos se segregaron por responsabilidad específica:

\begin{table}[h!]
\centering
\renewcommand{\arraystretch}{1.35}
\begin{tabularx}{\textwidth}{l l X}
\toprule
\textbf{Capa} & \textbf{Contrato / Interfaz} & \textbf{Segregación Aplicada} \\
\midrule
Dominio (.NET) & \texttt{IBusinessRepository} & Expone exclusivamente operaciones de negocio (\texttt{GetBySlugAsync}, \texttt{SlugExistsAsync}, \texttt{GetByIdAsync}, \texttt{UpdateAsync}) sin mezclar consultas de citas o usuarios. \\
Dominio (.NET) & \texttt{IUnitOfWork} & Expone únicamente el compromiso transaccional (\texttt{SaveChangesAsync}), separando la persistencia de las consultas de lectura (\texttt{IReadRepository<T>}). \\
UI (React Native) & \texttt{FormBrandIdentityCardProps} & Recibe únicamente \texttt{businessName}, \texttt{logoUrl}, \texttt{accentColor}, \texttt{secondaryColor} y \texttt{onChangeBranding}, sin depender del arreglo de preguntas \texttt{fields}. \\
\bottomrule
\end{tabularx}
\caption{Segregación de contratos e interfaces en el módulo de identidad de marca.}
\end{table}

% ============================================================================
\section{D --- Dependency Inversion Principle (Inversión de Dependencias)}
% ============================================================================

\begin{quote}
\textit{``Los módulos de alto nivel no deben depender de módulos de bajo nivel. Ambos deben depender de abstracciones.''}
\end{quote}

\begin{enumerate}[label=\textbf{\arabic*.}]
    \item \textbf{Desacoplamiento en Controladores REST (\texttt{BusinessesController.cs}):} El controlador de la capa \texttt{Slotify.API} no depende de \texttt{AppDbContext} ni de repositorios concretos de \texttt{Slotify.Infrastructure}; depende exclusivamente de la abstracción \texttt{IMediator}:
\begin{lstlisting}[language={[Sharp]C}]
[ApiController]
[Route("api/businesses")]
public class BusinessesController(IMediator mediator) : ControllerBase
{
    private readonly IMediator _mediator = mediator;
    // Envía comandos y consultas a través de la abstracción IMediator
}
\end{lstlisting}
    \item \textbf{Inyección de Abstracciones en Casos de Uso (\texttt{Slotify.Application}):} \texttt{UpdateBusinessBrandIdentityCommandHandler} recibe por constructor las interfaces \texttt{IBusinessRepository} e \texttt{IUnitOfWork} definidas en \texttt{Slotify.Domain}. La capa de aplicación desconoce por completo si el almacenamiento subyacente es PostgreSQL, Supabase o memoria volátil.
    \item \textbf{Abstracción de Almacenamiento Local en Frontend (\texttt{storage.ts} y \texttt{formBranding.ts}):} Los componentes visuales invocan las funciones de alto nivel \texttt{saveFormSchemaLocally} y \texttt{loadFormSchemaLocally}, las cuales delegan en el adaptador abstracto \texttt{storage} (que gestiona internamente \texttt{AsyncStorage}, \texttt{SecureStore} o memoria en caché sin exponer detalles nativos a la vista).
\end{enumerate}

% ============================================================================
\section{Matriz de Trazabilidad SOLID}
% ============================================================================

\begin{table}[h!]
\centering
\renewcommand{\arraystretch}{1.35}
\begin{tabularx}{\textwidth}{c l X}
\toprule
\textbf{Principio} & \textbf{Artefactos Clave} & \textbf{Garantía de Calidad Obtenida} \\
\midrule
\textbf{S} & \texttt{Business.cs}, \texttt{FormBrandIdentityCard.tsx}, \texttt{formBranding.ts} & Alta cohesión; separación estricta entre reglas de dominio, validación hexadecimal, UI de edición y visor público. \\
\textbf{O} & \texttt{ValidationBehavior.cs}, \texttt{CORPORATE\_COLOR\_PRESETS} & Nuevas paletas corporativas, logotipos y validadores se agregan sin modificar el código existente probado. \\
\textbf{L} & \texttt{IBusinessRepository}, \texttt{BookingFormConfigDto} & Sustitución transparente por \textit{Mocks} en pruebas unitarias y retrocompatibilidad con plantillas de sector. \\
\textbf{I} & \texttt{FormBrandIdentityCardProps}, \texttt{IReadRepository<T>} & Contratos reducidos y específicos; los componentes no dependen de propiedades que no consumen. \\
\textbf{D} & \texttt{IMediator}, \texttt{IBusinessRepository}, \texttt{IUnitOfWork} & Inyección de dependencias (DI) en todas las capas; cero acoplamiento directo a EF Core o detalles de infraestructura. \\
\bottomrule
\end{tabularx}
\caption{Resumen de cumplimiento de los principios SOLID en la historia de usuario.}
\end{table}

\end{document}

